% =========================================================
% Common Induction Mistakes
% =========================================================
\paragraph{Common Induction Mistakes.}

\begin{remark}[Mistake 1: Missing base case]
A proof that establishes the inductive step but not the base case proves
nothing. The step says: \emph{if} $P(n)$ holds for some $n$, \emph{then}
$P(n+1)$ holds. Without a base case, there is no $n$ for which $P(n)$
is established, so the chain never starts.

The classic illustration: the false theorem ``all horses are the same
colour.'' The base case for $n = 1$ is trivially true (one horse has
one colour). The inductive step has a subtle flaw in the overlap argument.
The lesson is that both parts are essential and independent. Always
verify the base case, even when it seems trivial. A trivial base case
takes one line; write the one line.
\end{remark}

\begin{remark}[Mistake 2: Circular inductive step]
The inductive step proves $P(n+1)$ using only $P(n)$ and previously
established facts. It is a fatal error to assume $P(n+1)$ at any point
in the proof of $P(n+1)$.

This error most often appears when a student ``derives'' both sides of
an equation and meets in the middle, inadvertently assuming what they
are trying to prove. The symptom: the final few lines of the step are
not deductions from $P(n)$ but rather manipulations of $P(n+1)$ itself.

The cure: always work in one direction. Start with the left side of
$P(n+1)$ (or whatever the goal is), apply the inductive hypothesis once,
and algebraically arrive at the right side. Never write the target as
an intermediate line and work toward it from both sides.
\end{remark}

\begin{remark}[Mistake 3: Unstated inductive hypothesis]
The most common error in practice is writing ``assume $P(n)$'' without
saying what $P(n)$ actually says. This is a sign that $P(n)$ has not
been written out as a full sentence.

The effect is that the inductive step becomes vague at the crucial
moment: the student needs to use the inductive hypothesis but cannot
cite it precisely, so they either paraphrase it loosely (which may
introduce errors) or implicitly use a stronger version than what was
stated.

Write $P(n)$ as a complete sentence at the top of the proof. Then
write ``assume $P(n)$: $[\text{full sentence}]$.'' This forces precision
and prevents the most common errors in the body of the proof.
\end{remark}

\begin{remark}[Mistake 4: Wrong base case]
Induction proves $P(n)$ for all $n \geq n_0$, where $n_0$ is whatever
value you use as the base case. If the claim is only true for $n \geq 2$,
the base case $n = 0$ will fail, and the proof is invalid.

Always check: what is the smallest $n$ for which the claim is true?
That is your base case. Do not default to $n = 0$ or $n = 1$ without
verifying that the claim holds there. For divisibility statements,
$n = 0$ often gives $0 = 0 \cdot k$, which works. For growth bounds,
$n = 1$ or $n = 2$ may be needed. For Fibonacci identities, specific
small values may require individual checking.
\end{remark}

\begin{remark}[Mistake 5: Using strong induction where weak suffices, and vice versa]
Strong induction is not automatically better than weak. A proof that
invokes the full strong hypothesis ``$P(k)$ for all $k < n$'' when
only $P(n-1)$ is actually used is unnecessarily complicated and obscures
the proof's logic.

Conversely, a proof that tries to use weak induction on a recursion
requiring two prior terms will stall at the step.

The diagnostic: write the step and check explicitly which prior values
you use. If only $P(n)$, use weak. If $P(k)$ for some $k < n$ that
is not $n-1$, use strong.
\end{remark}

\begin{remark}[Mistake 6: Forgetting vacuity in strong induction]
In strong induction, the case $n = 0$ has a vacuously true hypothesis
$\forall k < 0,\; P(k)$, since no natural number is less than $0$.
This means $P(0)$ must be proved \emph{directly}, not from the inductive
hypothesis. Many students write ``by the inductive hypothesis applied
to $n - 1$'' for $n = 0$, which refers to $n - 1 = -1$, a non-existent
value.

The base case in strong induction is always proved from scratch.
The strong hypothesis is only available for $n \geq 1$.
\end{remark}
